% Source: https://artofproblemsolving.com/community/c3h1917968_math_problems
% Stable identity: aops-topic-1917968-q03
The centre $O$ is equidistant from $(6,0)$ and $(8,0)$, so it lies on their perpendicular bisector $x=7$; write $O=(7,k)$. Equal distances to $(8,0)$ and $(3,2)$ give
\[
1^2+k^2=4^2+(k-2)^2\iff 1=16-4k+4\iff k=\dfrac{19}{4}.
\]

\begin{asy}
usepackage("amsmath");
size(6.5cm);
defaultpen(fontsize(10pt));
pair A=(6,0),B=(8,0),P=(3,2),O=(7,19/4); real r=sqrt(377)/4;
draw((-1,0)--(13,0),Arrow(6)); draw((0,-1)--(0,10.5),Arrow(6));
label("$x$",(13,0),E); label("$y$",(0,10.5),N);
draw(circle(O,r),blue+linewidth(1));
draw((7,-0.8)--(7,10),dashed);
draw(O--A,red+linewidth(0.8)); draw(O--P,red+linewidth(0.8));
dot(A);dot(B);dot(P);dot(O);
label("$(6,0)$",A,SW); label("$(8,0)$",B,SE); label("$(3,2)$",P,W);
label("$O\left(7,\dfrac{19}{4}\right)$",O,E); label("$r$",(O+A)/2,W,red); label("$x=7$",(7,10),N);
label("\mbox{}",(min(currentpicture,user=true)+max(currentpicture,user=true))/2,black); // keeps the text after the figure black when LaTeX typesets the labels inline
\end{asy}

Then $r^2=1+\left(\dfrac{19}{4}\right)^2=\dfrac{16+361}{16}=\dfrac{377}{16}$, so the radius is $\boxed{\dfrac{\sqrt{377}}{4}}\approx4.85$.
